% chapter4.tex — บทที่ 4 ผลการวิเคราะห์ข้อมูล
\chapter{ผลการวิเคราะห์ข้อมูล}

\section{สัญลักษณ์ที่ใช้ในการนำเสนอผลการวิเคราะห์ข้อมูล}
(ระบุสัญลักษณ์ทางสถิติที่ใช้ เช่น $\bar{x}$ แทนค่าเฉลี่ย, $S.D.$ แทนส่วนเบี่ยงเบนมาตรฐาน)

\section{ผลการวิเคราะห์ข้อมูล}
(นำเสนอผลการวิเคราะห์ข้อมูลเรียงตามวัตถุประสงค์/คำถามวิจัยแต่ละข้อ ประกอบตารางและการแปลผล)

\begin{longtable}{|p{4.5cm}|c|c|c|c|}
\caption{ตัวอย่างรูปแบบตารางแสดงผลการวิเคราะห์ข้อมูล} \label{tab:4-1}\\
\hline
\tblhead{รายการ} & \tblhead{จำนวน (n)} & \tblhead{ค่าเฉลี่ย ($\bar{x}$)} & \tblhead{S.D.} & \tblhead{ระดับ} \\
\hline
\endfirsthead
\multicolumn{5}{c}{\tablename\ \thetable\ (ต่อ)}\\
\hline
\tblhead{รายการ} & \tblhead{จำนวน (n)} & \tblhead{ค่าเฉลี่ย ($\bar{x}$)} & \tblhead{S.D.} & \tblhead{ระดับ} \\
\hline
\endhead
\hline
\endfoot
\hline
\endlastfoot
{[}รายการที่ 1{]} & 0 & 0.00 & 0.00 & {[}ระดับ{]} \\
\hline
{[}รายการที่ 2{]} & 0 & 0.00 & 0.00 & {[}ระดับ{]} \\
\hline
รวม & 0 & 0.00 & 0.00 & {[}ระดับ{]} \\
\hline
\end{longtable}
